% Round 15. Modular reduction; A2 is discharged by the final actual projection theorem.
\section{A fixed-section kernel reducing full Aeppli degree two}
\label{r15fs:section}

This section retains a modular hypothesis to isolate exactly
which actual projection degrees are needed. That hypothesis is
discharged by Theorem~\ref{r15a2:projection}, and the final
activation is recorded in Corollary~\ref{r15a2:full-HA22}.
The argument does not assume an output in degree $(2,2)$
or a base-change theorem for new currents.

\begin{definition}[The three-input actual projection hypothesis]
\label{r15fs:hypothesis}
Let $(\mathrm A_2)$ denote the following assertion. For every
smooth integral projective $K$-variety $B$, projection
$B\times\mathbf P^1\to B$ sends the proper current of
each ordinary actual input in the three degrees
\[
                         (2,2),\quad(3,2),\quad(2,3)
\tag{FX15.1}\label{r15fs:input-degrees}
\]
to a current having an ordinary actual representative in
the respective degrees $(1,1),(2,1),(1,2)$, with agreement
against every ordinary smooth compact mixed test.
The inputs are not assumed closed. The assertion may
equivalently be supplied locally over complete proper
fibers, with ordinary gluing by current injectivity.
No base-change statement is included in this hypothesis.
\end{definition}

\begin{proposition}[The fixed-section homotopy]
\label{r15fs:homotopy}
Assume $(\mathrm A_2)$. Let $S/K$ be smooth integral
projective, let $Y=S\times\mathbf P^1$, and let
$\alpha\in B^{2,2}(Y)$ satisfy $D\alpha=0$.
Fix $z_0\in\mathbf P^1(K)$ and put
$\alpha_0=(\mathrm{id}_S,z_0)^*\alpha$.
There are ordinary actual forms
\[
 R_\alpha\in B^{1,1}(Y),\quad
 G_\alpha^0\in B^{1,1}(Y),\quad
 P_\alpha^0\in B^{2,1}(Y),\quad
 Q_\alpha^0\in B^{1,2}(Y)
\]
such that $DR_\alpha=0$ and
\[
 \alpha=h_y\wedge R_\alpha+\pi^*\alpha_0
               +DG_\alpha^0+d''P_\alpha^0-d'Q_\alpha^0.
\tag{FX15.2}\label{r15fs:formula}
\]
Here $h_y$ is a fixed smooth separately closed curvature
of degree one on the projective-line factor. The current
of $R_\alpha$ is $(p_y)_*[p_z^*\alpha]$; it need not
be identified with the pullback of a new base output in
order to use the formula.
\end{proposition}

\begin{proof}
On $T=S\times\mathbf P^1_z\times\mathbf P^1_y$, take
the already fixed diagonal kernel $g_\Delta$, and on
$\mathbf P^1_z$ take the scalar point Green function
$g_0$ for $z_0$ with the same smooth reference curvature.
Both are ordinary PWS away from their indicated divisors,
and their scalar Green identities give
\[
 Dg_\Delta=h_z+h_y-\delta_\Delta,
 \qquad Dg_0=h_z-\delta_{z_0}.
\]
Thus the PWS scalar $g^0=g_\Delta-g_0(z)$ satisfies
\[
                     Dg^0=h_y+\delta_{z_0}-\delta_\Delta.
\tag{FX15.3}\label{r15fs:kernel}
\]
The extra divisor has multiplicity and native degree one
because $z_0$ is $K$-rational.

We check ordinary excision for each transgression.
Write $n=\dim Y$. Near the diagonal the actual
immersion-germ theorem identifies $p_z^*\alpha$ with
$p_y^*\alpha$. Near the divisor $z=z_0$ it identifies
$p_z^*\alpha$ with $p_y^*\pi^*\alpha_0$. For tests
complementary to the three proposed outputs $G^0,P^0,Q^0$,
the products of these target degree-$(2,2)$ forms with
the tests have degrees respectively
\[
                    (n+1,n+1),\quad(n,n+1),\quad(n+1,n).
\tag{FX15.4}\label{r15fs:overdegree}
\]
They vanish by dimension. Hence all paired singular
integrands vanish on whole ordinary neighborhoods of
both divisors. Their intersection introduces no extra
condition: the union of the two zero neighborhoods
already suffices. Compactness permits one smooth cutoff
vanishing near that union and equal to one outside a
slightly larger zero neighborhood.

The cutoff inputs $g^0p_z^*\alpha$,
$d'g^0\wedge p_z^*\alpha$, and
$d''g^0\wedge p_z^*\alpha$, extended by zero, are
ordinary actual in exactly the degrees
\eqref{r15fs:input-degrees}. Hypothesis $(\mathrm A_2)$
gives the three ordinary representatives. Removing the
cutoff in any prescribed pairing changes only a zero
integrand. In particular no improper integral along a
root tail is used. The point Green function is regular
at infinity when $z_0$ is the chosen finite point;
the same argument is coordinate independent.

Apply $(\mathrm A_2)$ also to the global input
$p_z^*\alpha$ to obtain $R_\alpha$. Proper current
differentiation and injectivity of the actual-current
map imply $DR_\alpha=0$.

Finally apply the scalar Green identity
\eqref{r15fs:kernel} to $p_z^*\alpha\wedge p_y^*\eta$.
The $h_y$ term is $h_yR_\alpha$ by the ordinary smooth
target projection formula. Restriction to the divisor
$z=z_0$ gives $\pi^*\alpha_0$ against every smooth
mixed test, since that divisor maps isomorphically to
$Y$. The diagonal term gives $\alpha$. Expanding
$D(\alpha\wedge\eta)$ and using $D\alpha=0$, the
same two integrations by parts as in
\eqref{r15s:homotopy} give the signs
$+DG^0+d''P^0-d'Q^0$. All products near the divisors
have just been excised on ordinary zero neighborhoods,
so this calculation does not multiply an undefined
singular current by an actual form. It proves
\eqref{r15fs:formula} as currents on all smooth tests,
and hence as actual forms by current injectivity.
\end{proof}

\begin{corollary}[The full degree-two formula for a trivial projective line]
\label{r15fs:product-formula}
Assume $(\mathrm A_2)$. For every smooth integral projective
$S/K$, with $\pi:Y=S\times\mathbf P^1\to S$, one has
\[
 \HA^{2,2}(Y)=\pi^*\HA^{2,2}(S)
                       \oplus[h]\wedge\pi^*\HA^{1,1}(S).
\tag{FX15.9}\label{r15fs:full-product}
\]
The groups are the full actual groups. This statement
does not assert the analogous formula for a nontrivial
projective-line bundle without a supplied section kernel.
\end{corollary}

\begin{proof}
The homotopy \eqref{r15fs:formula} and the already proved
full $(1,1)$ formula give
$[R_\alpha]=\pi^*[b]+c[h]$ for a $D$-closed actual
base class $[b]$. Since $h^2=0$ strictly and multiplication
by $h$ preserves actual Aeppli boundaries, they give
\[
                 [\alpha]=\pi^*[\alpha_0]+[h]\pi^*[b].
\]
This proves spanning.

For directness, pull a putative zero relation
$\pi^*[a]+[h]\pi^*[b]=0$ back by the fixed algebraic
section. The pullback of $h$ to that section is zero,
so $[a]=0$. Projection of an actual Aeppli boundary
$d'\mu+d''\nu$ in degree $(2,2)$ is again an actual
Aeppli boundary: the first-strip outputs of $\mu$ and
$\nu$ have degrees $(0,1),(1,0)$, and proper current
differentiation identifies their derivatives with the
projected current. Thus the operator on degree-$(2,2)$
inputs supplied by $(\mathrm A_2)$ induces a map on the actual Aeppli
quotient. This map sends $[h]\pi^*[b]$ to $[b]$,
by frozen formula (R10.1) applied to every actual
top test $b\wedge\eta$ on the base. It follows that
$[b]=0$, proving the direct sum. No current base-change
formula is used.
\end{proof}

\begin{corollary}[Conditional full degree-two product and projective-space groups]
\label{r15fs:projective-space}
Assume $(\mathrm A_2)$ and the already proved full actual
$(1,1)$ projective-bundle formula. Then
\[
 H_{\mathrm A,\mathrm{act}}^{2,2}((\mathbf P^1)^m)
       =\bigoplus_{1\leq i<j\leq m}\mathbf R[h_i h_j].
\tag{FX15.5}\label{r15fs:product}
\]
For $m\geq2$ one consequently has
\[
 H_{\mathrm A,\mathrm{act}}^{2,2}(\mathbf P^m_K)
                              =\mathbf R[h^2].
\tag{FX15.6}\label{r15fs:space}
\]
For $m<2$ the group is zero by degree. These are full
actual groups, conditional on $(\mathrm A_2)$; no
claim here discharges that hypothesis. The same statements
hold on the corresponding spaces over any finite $L/K$,
viewed in the native $K$ theory; their top periods are
then $[L:K]$ instead of one.
\end{corollary}

\begin{proof}
Apply the preceding proposition with
$S=(\mathbf P^1)^{m-1}$. The already proved actual
$(1,1)$ formula writes
\[
            [R_\alpha]=\pi^*[\beta]+c[h_m],
                 \qquad [\beta]\in\HA^{1,1}(S).
\]
Multiplication by the smooth separately closed $h_m$
preserves actual Aeppli boundaries. Moreover $h_m^2=0$
as a form, since it is pulled from a curve. Therefore
\eqref{r15fs:formula} gives
\[
           [\alpha]=\pi^*[\alpha_0]+h_m\pi^*[\beta].
\tag{FX15.7}\label{r15fs:induction}
\]
This step does not require a base-change identity for
$R_\alpha$ or an operator projecting a $(3,3)$ input.
Induction, starting with the dimension-one zero group,
gives the spanning in \eqref{r15fs:product}.
Pairing with $\prod_{k\notin\{i,j\}}h_k$ gives the
coefficient of $h_i h_j$: every other monomial repeats
some $h_k$, and $\int\prod_kh_k=1$ in the native
$K$ theory. Actual Stokes kills every actual boundary
against these separately closed tests. This proves
directness in the full actual quotient.

Use the finite symmetric quotient
$q:(\mathbf P^1)^m\to\mathbf P^m$, of native degree
$m!$. Its pullback is injective on full actual Aeppli
degree $(2,2)$ by frozen Corollary
\texttt{r14h:other-cohomology}; the boundary primitives
have degrees $(1,2),(2,1)$ and lie in the known finite
actual trace strip. The permutation invariance of
$q^*[\alpha]$ makes all coefficients in
\eqref{r15fs:product} equal to some $c$.
The line-bundle identity $q^*\mathcal O(1)=
\bigotimes_i\mathcal O_i(1)$ and a smooth metric ratio
give $q^*h=\sum_i h_i+Du$ for a global smooth scalar.
The difference of the squares is $D$-exact because
the curvature factors are separately closed. Hence
\[
                  q^*[h^2]=2\sum_{i<j}[h_i h_j].
\]
Injectivity now gives $[\alpha]=(c/2)[h^2]$.
The class $[h^2]$ is nonzero by its native pairing
with $h^{m-2}$, whose value is one.
For spaces over finite $L/K$, all bases in the argument
are still smooth integral projective $K$-varieties and
are included in $(\mathrm A_2)$. Every coefficient
period is divided by $[L:K]$, while the symmetric
quotient still has degree $m!$ because its source and
target have the same constant field. No new actual
finite trace in degree $(2,2)$ is used.
\end{proof}

\begin{corollary}[Conditional Bott--Chern descent through point blowups]
\label{r15fs:point-blowups}
Assume $(\mathrm A_2)$. Let $X/K$ be smooth integral
projective and let $q:Y\to X$ be its blowup at finitely
many closed points. Then
\[
 q^*:H_{\mathrm{BC,act}}^{3,3}(X)
                      \longrightarrow H_{\mathrm{BC,act}}^{3,3}(Y)
\tag{FX15.8}\label{r15fs:point-BC}
\]
is injective. The same assertion holds for a finite
composition of such point blowups. The upstairs
primitive may be an arbitrary actual $(2,2)$ form;
no PWS primitive or exceptional representability
certificate must be supplied separately.
\end{corollary}

\begin{proof}
For one closed point $p$, its exceptional divisor is
$E=\mathbf P^{n-1}_{\kappa(p)}$, where $n=\dim X$.
Suppose a separately closed actual $(3,3)$ form
$\Omega$ has $q^*\Omega=D\beta$ with
$\beta\in B^{2,2}(Y)$ arbitrary actual. The restriction
$i^*\beta$ is $D$-closed because the center is a point
and the restriction of $\Omega$ to it is zero.
The preceding conditional computation, including its
native constant-field version, places its full actual
Aeppli class in $\mathbf R[h^2]$ (or zero if the
exceptional dimension is below two).

The square of the global separately closed model
curvature of $\mathcal O_Y(-E)$ restricts to this
hyperplane-square class. A change to a smooth metric
changes the square by an actual $D$ boundary, as in
Theorem~\ref{r15pc:point-certificate}. Thus the
exceptional class lies in the restriction image from
$Y$. The all-bidegree connecting sequence
of Proposition~\ref{r15m:mixed-connecting}, in degree
$(2,2)$, now gives $\Omega=D\Psi$ on $X$ for an
ordinary actual $\Psi\in B^{2,2}(X)$.

For disjoint centers this argument applies componentwise.
For a finite sequence, the primitive returned at one
step is actual, which is exactly the input class allowed
at the next step; hence the argument iterates. It does
not prove descent for positive-dimensional centers,
whose exceptional groups have a base, and it does not
prove finite Bott--Chern injectivity in degree $(3,3)$,
which would require a finite actual trace of a general
$(2,2)$ primitive.
\end{proof}
